\scanpage{97}
\mspage{45}

On pose, si $Z$ est une partie de $X$,
\[(99)\qquad\mathcal V_{i_0,\ldots,i_n}^{Z}=\mathcal V_{i_0\ldots i_n}\cap Z\qquad(\hookrightarrow\mathcal V_{i_0\ldots i_n})\]
si $Z$ est \deleted{fermé} une partie \unsure{loc\textsuperscript{t}|loc!} fermée qui contient \inserted{$X_{i_n}^*$ (on sait que $\unsure{X_{i_n}^*\cap\mathcal V_{i_0\ldots i_n}}\overset{(\overset h\sim)}{\hookrightarrow}\mathcal V_{i_0\ldots i_n}$)}
\deleted{\[X_{i_n}-\dot X_{i_0}\qquad(\mathrm{NB}\ \dot X_{i_0}\overset{\mathrm{déf}}=X_{i_0}\setminus X_{i_0}^*=\bigcup_{i<i_0}X_i=\bigcup_{i<i_0}X_i^*)\]}
alors on a
\[(100)\qquad\mathcal V_{i_0\ldots i_n}^{Z}\overset{(\overset h\sim)}{\hookrightarrow}\mathcal V_{i_0\ldots i_n}\qquad(\text{équiv. d'homotopie})\]
\margin{?\quad\uline{à vérifier!}}
(on suppose $X$ métrisable\ldots)

On va généraliser la définition des $\mathcal V_i$, $\mathcal V_{i_0\ldots i_n}$ en des $\mathcal V_J$, $\mathcal V_{J_0\ldots J_n}$, où $J$ est une bande de $I$, resp. $J_0\prec\cdots\prec J_n$ est un ``drapeau de bandes'' de $I$. On définit
\[
(101)\qquad\begin{aligned}
\mathcal V_J&=\mathcal V(X_J^*,X)\qquad(\text{provoisinage tubulaire})\\
&=\bigcup_{j\in J}\mathcal V_j=\underset{d\in\mathrm{Drap}(J)}{\lim_{\longrightarrow}}\mathcal V_d
\end{aligned}
\]
(NB la définition et les deux formules \unsure{marchent} encore pour $J$ partie quelconque de $I$)
\[(102)\qquad\mathcal V_J^Z=\mathcal V_J\cap Z\]
On a donc
\[(103)\qquad X_J^*\overset{(\overset h\sim)}{\hookrightarrow}\mathcal V_J\qquad(\text{équiv. d'hom.})\]
\[(104)\]
\diagram{\begin{tikzcd}[ampersand replacement=\&,column sep=large,row sep=small]
X_J^*\arrow[r,hook,"(\overset h\sim)"]\&\mathcal V_J^Z\arrow[r,hook,"(\overset h\sim)"]\arrow[d,equal]\&\mathcal V_J\\
\&\mathcal V(X_J^*,Z)\&
\end{tikzcd}}
(équiv. d'hom) \uuline{si} $Z\supset X_J$
\margin{96}
